% c00 原书前言（Second Edition p.iii；First Edition p.v–vi）PDF p.004–006
\section*{第二版序}
\addcontentsline{toc}{section}{第二版序}

本书是 2006 年由清华大学出版社出版的、题名相同的书的第二版。与第一版一样，我们的目标是以最少的预备知识发展概形上凝聚层的上同调理论。我们覆盖 Grothendieck《Éléments de Géométrie Algébrique》(EGA) I–III 的主要材料，并假定读者熟悉 [Atiyah–Macdonald] 的第 1–8 章。我们改正了第一版的错误，改进了表述，并增加了习题。感谢赵天夫，使本书得以在高等教育出版社出版。

\vspace{1.5em}
\begin{flushright}
扶磊\\
丘成桐数学科学中心
\end{flushright}

\clearpage
\section*{第一版序}
\addcontentsline{toc}{section}{第一版序}

本书研究概形上凝聚层的上同调。这个题目的经典教材是 [Hartshorne]，但它有一些局限。例如，许多对固有态射成立的重要定理，在 [Hartshorne] 中只对射影态射证明。此外，[Hartshorne] 几乎不涉及谱序列技术。本书的主要目的就是填补这一空白。当然，本书中的一切都能在 Grothendieck 的 [EGA] 中找到。我希望本书能把隐藏在卷帙浩繁的 [EGA] 中的 Grothendieck 的奇妙思想讲得更清楚。为使本书自足，我们收录了许多在 [Hartshorne] 中处理得很好的材料。因此，本书与 [Hartshorne] 有一些重叠，尤其是在 \S1.2 与 \S1.4。

本书粗线条地覆盖 [EGA] I、II、III 的主要材料。数年前，我应邀在中国科学院晨兴数学中心 (MCM) 作一系列关于 $\ell$-adic 上同调理论的报告。为准备这些报告，我写了一本 étale 上同调理论的书，后来由 World Scientific 出版。它覆盖 [SGA] 1、4、$4\frac{1}{2}$、5、7 的主要材料。本书与 [Matsumura] 合起来，覆盖了阅读这本书所需的全部预备知识。所以我以为，本书能为学习 étale 上同调理论提供充分的准备。

我只假定读者熟悉 [Atiyah–Macdonald] 的第 1–8 章。本书用到的其他代数结果，或在本书中证明，或读者不难自行证明。我收录一些代数材料，是因为我相信：学代数的最好方式是与几何一同去学，这样才能对抽象的代数概念获得几何直观。

本书绝非代数几何的完备论著。对于如何用本书中由上同调方法得到的结果去研究代数簇的几何，本书一字未提；本书也不包含 Serre 对偶。这些题目请读者查阅 [Hartshorne] 及其参考文献。

我衷心感谢冯克勤教授。自从他知道我的存在以来，他从未停止对我的鼓励。他邀请我在 MCM 讲授 $\ell$-adic 上同调理论与其他代数几何专题，我的一些讲义构成本书的一部分。冯教授还使本书的出版成为可能。没有他的帮助，本书永远不会问世。

在本书的写作过程中，我得到了求是科技基金会、973 项目、IHES 与 MCM 的支持。

\vspace{1.5em}
\begin{flushright}
扶磊\\
陈省身数学研究所
\end{flushright}
